\section{Introduction}
\label{sec:introduction}

Segmentation of the retinal vessels in fundus images gives quantitative markers for the diagnosis and follow-up of diseases such as diabetic retinopathy and glaucoma~\cite{jin2022fives, moccia2018blood}. Thin vessels are a particular challenge in this task: they are a few pixels wide, have low contrast, and weigh little in overlap metrics such as Dice, in which a wide vessel counts more pixels than a thin vessel of the same length.

Convolutional neural networks (CNNs) have dominated vessel segmentation since the U-Net~\cite{ronneberger2015u, siddique2021u}, and a minimal U-Net with far fewer parameters than the best published models, carefully trained, closely approximates the best published results on retinal benchmarks~\cite{galdran2020littlewnetcouldstateoftheart}. Vision Transformers (ViTs)~\cite{dosovitskiy2020image} and the foundation models built on them, such as the Segment Anything Model~\cite{Kirillov_2023_ICCV} and its medical version \medsam{}~\cite{ma2024segment}, offer large pre-trained encoders instead. A plain ViT, however, splits the image into fixed-size patches, $16 \times 16$ pixels in \medsam{}~\cite{ma2024segment}, and encodes each patch as one token. Transformers that encode such patches without convolutional features have been reported to lose the low-level detail needed for precise localization~\cite{chen2021transunet}, which is a strong reason to expect a plain ViT to miss thin vessels. It is not the only one: plain ViTs lack the locality bias of convolutions and reach their accuracy only when pre-trained on large datasets~\cite{dosovitskiy2020image}, while vessel datasets have tens to hundreds of annotated images.

Testing this expectation requires measuring the thin vessels separately from the overall score and comparing the architectures with a training protocol that aims at the same conditions for all of them. We ask three questions:
\begin{enumerate}
    \item Can a plain ViT (\medsam{}) reach the segmentation performance of CNNs on thin vessels?
    \item How do the plain and the hybrid Transformer compare with CNNs in thin-vessel performance relative to their cost, in parameters and operations per image, and can a compact plain ViT keep the thin-vessel performance at the cost of a CNN?
    \item How much segmentation performance on thin vessels does foundation pre-training add to transformer-based models, and how does this gain depend on the number of training images?
\end{enumerate}

We trained the plain ViT \medsam{}, the hybrid \litemedsam{} (convolutions and attention) and three CNNs of different sizes, \convnext{}~\cite{liu2022convnet}, \convunext{}~\cite{han2022convunext} and \lwnet{}~\cite{galdran2020littlewnetcouldstateoftheart}, on three retinal datasets, DRIVE~\cite{Staal2004}, HRF~\cite{budai2013robust} and FIVES~\cite{jin2022fives}, and on VessMap~\cite{viana2025new}, a microscopy dataset that serves as a check in a second imaging domain. The training protocol aimed at the same conditions for every model: a search of 100 trials, the same loss, and the same test images and metrics (Appendix~\ref{app:training} gives where validation differed). The two Transformers were also trained from scratch and as compact models. Besides Dice, we measured the centerline recall by vessel radius and the centerline recall and precision of thin vessels, and compared the models with paired tests over the test images.

Fine-tuned \medsam{}, evaluated at the scale it was trained on, had a centerline recall of thin vessels not significantly different from that of the best CNN, the CNN with the highest mean for each dataset and metric, on all four datasets, and a significantly lower centerline precision of thin vessels on DRIVE and VessMap. This depended on pre-training and, on the small DRIVE and VessMap images, on evaluating them upscaled as in training. A compact plain ViT trained from scratch did not keep this performance. On HRF and FIVES, fine-tuned \medsam{} also had a significantly higher Dice than the best CNN.
